% sec1_2.tex — Section 1.2: The Algebra of Least Squares

This section describes the computational procedure for obtaining the OLS estimate,
$\vec{b}$, of the unknown coefficient vector $\bm{\beta}$ and introduces a few concepts that derive
from~$\vec{b}$.

%% ─── OLS Minimizes the Sum of Squared Residuals ──────────────────────────────
\subsection*{OLS Minimizes the Sum of Squared Residuals}

Although we do not observe the error term, we can calculate the value implied by
a hypothetical value, $\widetilde{\bm{\beta}}$, of $\bm{\beta}$ as
\[
  y_i - \vec{x}_i' \widetilde{\bm{\beta}}.
\]
This is called the \textbf{residual} for observation~$i$.  From this, form the
\textbf{sum of squared residuals (SSR)}:
\[
  SSR(\widetilde{\bm{\beta}})
    \equiv \sum_{i=1}^{n} (y_i - \vec{x}_i' \widetilde{\bm{\beta}})^2
    = (\vec{y} - \vec{X}\widetilde{\bm{\beta}})'
      (\vec{y} - \vec{X}\widetilde{\bm{\beta}}).
\]
This sum is also called the \textbf{error sum of squares (ESS)} or the
\textbf{residual sum of squares (RSS)}.  It is a function of
$\widetilde{\bm{\beta}}$ because the residual depends on it.  The \textbf{OLS
estimate}, $\vec{b}$, of $\bm{\beta}$ is the $\widetilde{\bm{\beta}}$ that
minimizes this function:
\begin{equation}
  \vec{b} \equiv \argmin_{\widetilde{\bm{\beta}}} SSR(\widetilde{\bm{\beta}}).
  \label{eq:1.2.1}
\end{equation}

The relationship among $\bm{\beta}$ (the unknown coefficient vector), $\vec{b}$
(the OLS estimate of it), and $\widetilde{\bm{\beta}}$ (a hypothetical value of
$\bm{\beta}$) is illustrated in Figure~\ref{fig:1.1} for $K = 1$.  Because
$SSR(\widetilde{\bm{\beta}})$ is quadratic in $\widetilde{\bm{\beta}}$, its graph
has the U~shape.  The value of $\widetilde{\bm{\beta}}$ corresponding to the bottom
is~$\vec{b}$, the OLS estimate.  Since it depends on the sample
$(\vec{y}, \vec{X})$, the OLS estimate $\vec{b}$ is in general different from the
true value $\bm{\beta}$; if $\vec{b}$ equals $\bm{\beta}$, it is by sheer
accident.

By having squared residuals in the objective function, this method imposes a heavy
penalty on large residuals; the OLS estimate is chosen to prevent large residuals
for a few observations at the expense of tolerating relatively small residuals for
many other observations.  We will see in the next section that this particular
criterion brings about some desirable properties for the estimate.

\begin{figure}[htbp]
\centering
\includegraphics[width=0.9\textwidth,trim=50 400 50 400,clip]{figures/fig_1_1.png}
\caption{Hypothetical, True, and Estimated Values}
\label{fig:1.1}
\end{figure}

%% ─── Normal Equations ─────────────────────────────────────────────────────────
\subsection*{Normal Equations}

A sure-fire way of solving the minimization problem is to derive the first-order
conditions by setting the partial derivatives equal to zero.  To this end we seek a
$K$-dimensional vector of partial derivatives,
$\partial SSR(\widetilde{\bm{\beta}})/\partial \widetilde{\bm{\beta}}$.\footnote{If
$h \colon \mathbb{R}^K \to \mathbb{R}$ is a scalar-valued function of a
$K$-dimensional vector $\vec{x}$, the derivative of~$h$ with respect to $\vec{x}$
is a $K$-dimensional vector whose $k$-th element is $\partial h(\vec{x})/\partial x_k$
where $x_k$ is the $k$-th element of~$\vec{x}$.  (This $K$-dimensional vector is
called the \textbf{gradient}.)  Here, the $\vec{x}$ is $\widetilde{\bm{\beta}}$
and the function $h(\vec{x})$ is $SSR(\widetilde{\bm{\beta}})$.}
The task is facilitated by writing $SSR(\widetilde{\bm{\beta}})$ as
\begin{align}
  SSR(\widetilde{\bm{\beta}})
    &= (\vec{y} - \vec{X}\widetilde{\bm{\beta}})'
       (\vec{y} - \vec{X}\widetilde{\bm{\beta}})
    &&\text{(since the $i$-th element of
       $\vec{y} - \vec{X}\widetilde{\bm{\beta}}$ is
       $y_i - \vec{x}_i'\widetilde{\bm{\beta}}$)} \notag \\
    &= (\vec{y}' - \widetilde{\bm{\beta}}'\vec{X}')
       (\vec{y} - \vec{X}\widetilde{\bm{\beta}})
    &&\text{(since
       $(\vec{X}\widetilde{\bm{\beta}})' = \widetilde{\bm{\beta}}'\vec{X}'$)}
       \notag \\
    &= \vec{y}'\vec{y}
       - \widetilde{\bm{\beta}}'\vec{X}'\vec{y}
       - \vec{y}'\vec{X}\widetilde{\bm{\beta}}
       + \widetilde{\bm{\beta}}'\vec{X}'\vec{X}\widetilde{\bm{\beta}}
       \notag \\
    &= \vec{y}'\vec{y}
       - 2\vec{y}'\vec{X}\widetilde{\bm{\beta}}
       + \widetilde{\bm{\beta}}'\vec{X}'\vec{X}\widetilde{\bm{\beta}}
       \notag \\
    &\qquad\text{(since the scalar
       $\widetilde{\bm{\beta}}'\vec{X}'\vec{y}$ equals its transpose
       $\vec{y}'\vec{X}\widetilde{\bm{\beta}}$)} \notag \\
    &\equiv \vec{y}'\vec{y}
       - 2\vec{a}'\widetilde{\bm{\beta}}
       + \widetilde{\bm{\beta}}'\vec{A}\widetilde{\bm{\beta}}
    \quad\text{with } \vec{a} \equiv \vec{X}'\vec{y}
    \text{ and } \vec{A} \equiv \vec{X}'\vec{X}.
  \label{eq:1.2.2}
\end{align}

The term $\vec{y}'\vec{y}$ does not depend on $\widetilde{\bm{\beta}}$ and so can
be ignored in the differentiation of $SSR(\widetilde{\bm{\beta}})$.  Recalling from
matrix algebra that
\[
  \frac{\partial(\vec{a}'\widetilde{\bm{\beta}})}
       {\partial \widetilde{\bm{\beta}}} = \vec{a}
  \quad\text{and}\quad
  \frac{\partial(\widetilde{\bm{\beta}}'\vec{A}\widetilde{\bm{\beta}})}
       {\partial \widetilde{\bm{\beta}}}
    = 2\vec{A}\widetilde{\bm{\beta}}
  \quad\text{for $\vec{A}$ symmetric},
\]
the $K$-dimensional vector of partial derivatives is
\[
  \frac{\partial SSR(\widetilde{\bm{\beta}})}
       {\partial \widetilde{\bm{\beta}}}
    = -2\vec{a} + 2\vec{A}\widetilde{\bm{\beta}}.
\]

The first-order conditions are obtained by setting this equal to zero.  Recalling
from \eqref{eq:1.2.2} that $\vec{a}$ here is $\vec{X}'\vec{y}$ and $\vec{A}$ is
$\vec{X}'\vec{X}$ and rearranging, we can write the first-order conditions as
\begin{equation}
  \underset{(K \times K)}{\vec{X}'\vec{X}} \;
  \underset{(K \times 1)}{\vec{b}} = \vec{X}'\vec{y}.
  \label{eq:1.2.3}
\end{equation}

Here, we have replaced $\widetilde{\bm{\beta}}$ by $\vec{b}$ because the OLS
estimate $\vec{b}$ is the $\widetilde{\bm{\beta}}$ that satisfies the first-order
conditions.  These $K$ equations are called the \textbf{normal equations}.

The vector of residuals evaluated at $\widetilde{\bm{\beta}} = \vec{b}$,
\begin{equation}
  \underset{(n \times 1)}{\vec{e}} \equiv \vec{y} - \vec{X}\vec{b},
  \label{eq:1.2.4}
\end{equation}
is called the vector of \textbf{OLS residuals}.  Its $i$-th element is
$e_i \equiv y_i - \vec{x}_i'\vec{b}$.  Rearranging \eqref{eq:1.2.3} gives
\begin{equation}
  \vec{X}'(\vec{y} - \vec{X}\vec{b}) = \vec{0}
  \quad\text{or}\quad
  \vec{X}'\vec{e} = \vec{0}
  \quad\text{or}\quad
  \frac{1}{n}\sum_{i=1}^{n} \vec{x}_i \cdot e_i = \vec{0}
  \quad\text{or}\quad
  \frac{1}{n}\sum_{i=1}^{n} \vec{x}_i \cdot (y_i - \vec{x}_i'\vec{b}) = 0,
  \tag{1.2.3\ensuremath{'}}
  \label{eq:1.2.3prime}
\end{equation}
which shows that the normal equations can be interpreted as the sample analogue of
the orthogonality conditions $\E(\vec{x}_i \cdot \varepsilon_i) = \vec{0}$.  This
point will be pursued more fully in subsequent chapters.

To be sure, the first-order conditions are just a necessary condition for
minimization, and we have to check the second-order condition to make sure that
$\vec{b}$ achieves the minimum, not the maximum.  Those who are familiar with the
Hessian of a function of several variables\footnote{The \textbf{Hessian} of
$h(\vec{x})$ is a square matrix whose $(k, \ell)$ element is
$\partial^2 h(\vec{x})/\partial x_k \,\partial x_\ell$.} can immediately
recognize that the second-order condition is satisfied because (as noted below)
$\vec{X}'\vec{X}$ is positive definite.  There is, however, a more direct way to
show that $\vec{b}$ indeed achieves the minimum.  It utilizes the
``add-and-subtract'' strategy, which is effective when the objective function is
quadratic, as here.  Application of the strategy to the algebra of least squares
is left to you as an analytical exercise.

%% ─── Two Expressions for the OLS Estimator ───────────────────────────────────
\subsection*{Two Expressions for the OLS Estimator}

Thus, we have obtained a system of $K$ linear simultaneous equations in $K$
unknowns in~$\vec{b}$.  By Assumption~1.3 (no multicollinearity), the coefficient
matrix $\vec{X}'\vec{X}$ is positive definite (see review question~1 below for a
proof) and hence nonsingular.  So the normal equations can be solved uniquely for
$\vec{b}$ by premultiplying both sides of \eqref{eq:1.2.3} by
$(\vec{X}'\vec{X})^{-1}$:
\begin{equation}
  \vec{b} = (\vec{X}'\vec{X})^{-1}\vec{X}'\vec{y}.
  \label{eq:1.2.5}
\end{equation}

Viewed as a function of the sample $(\vec{y}, \vec{X})$, \eqref{eq:1.2.5} is
sometimes called the \textbf{OLS estimator}.  For any given sample
$(\vec{y}, \vec{X})$, the value of this function is the \textbf{OLS estimate}.  In
this book, as in most other textbooks, the two terms will be used almost
interchangeably.

Since $(\vec{X}'\vec{X})^{-1}\vec{X}'\vec{y}
= (\vec{X}'\vec{X}/n)^{-1}\vec{X}'\vec{y}/n$, the OLS estimator can also be
rewritten as
\begin{equation}
  \vec{b} = \vec{S}_{\vec{xx}}^{-1}\,\vec{s}_{\vec{xy}},
  \tag{1.2.5\ensuremath{'}}
  \label{eq:1.2.5prime}
\end{equation}
where
\begin{align}
  \vec{S}_{\vec{xx}}
    &= \frac{1}{n}\vec{X}'\vec{X}
     = \frac{1}{n}\sum_{i=1}^{n} \vec{x}_i \vec{x}_i'
    \quad\text{(sample average of $\vec{x}_i\vec{x}_i'$)},
  \label{eq:1.2.6a} \\
  \vec{s}_{\vec{xy}}
    &= \frac{1}{n}\vec{X}'\vec{y}
     = \frac{1}{n}\sum_{i=1}^{n} \vec{x}_i \cdot y_i
    \quad\text{(sample average of $\vec{x}_i \cdot y_i$)}.
  \label{eq:1.2.6b}
\end{align}

The data matrix form \eqref{eq:1.2.5} is more convenient for developing the
finite-sample results, while the sample average form \eqref{eq:1.2.5prime} is the
form to be utilized for large-sample theory.

%% ─── More Concepts and Algebra ───────────────────────────────────────────────
\subsection*{More Concepts and Algebra}

Having derived the OLS estimator of the coefficient vector, we can define a few
related concepts.

\begin{itemize}

\item The \textbf{fitted value} for observation~$i$ is defined as
  $\hat{y}_i \equiv \vec{x}_i'\vec{b}$.  The vector of fitted value, $\hat{\vec{y}}$,
  equals $\vec{X}\vec{b}$.  Thus, the vector of OLS residuals can be written as
  $\vec{e} = \vec{y} - \hat{\vec{y}}$.

\item The \textbf{projection matrix}~$\vec{P}$ and the \textbf{annihilator}~$\vec{M}$
  are defined as
  \begin{align}
    \underset{(n \times n)}{\vec{P}}
      &\equiv \vec{X}(\vec{X}'\vec{X})^{-1}\vec{X}',
    \label{eq:1.2.7} \\
    \underset{(n \times n)}{\vec{M}}
      &\equiv \vec{I}_n - \vec{P}.
    \label{eq:1.2.8}
  \end{align}
  They have the following nifty properties (proving them is a review question):
  \begin{align}
    &\text{Both $\vec{P}$ and $\vec{M}$ are symmetric and
      idempotent,}\footnote{A square matrix $\vec{A}$ is said to be
      \textbf{idempotent} if $\vec{A} = \vec{A}^2$.}
    \label{eq:1.2.9} \\
    \vec{P}\vec{X} &= \vec{X}
      \quad\text{(hence the term \textbf{projection} matrix)},
    \label{eq:1.2.10} \\
    \vec{M}\vec{X} &= \vec{0}
      \quad\text{(hence the term \textbf{annihilator})}.
    \label{eq:1.2.11}
  \end{align}

  Since $\vec{e}$ is the residual vector at
  $\widetilde{\bm{\beta}} = \vec{b}$, the sum of squared OLS residuals, $SSR$,
  equals $\vec{e}'\vec{e}$.  It can further be written as
  \begin{equation}
    SSR = \vec{e}'\vec{e} = \bm{\varepsilon}'\vec{M}\bm{\varepsilon}.
    \label{eq:1.2.12}
  \end{equation}
  (Proving this is a review question.)  This expression, relating $SSR$ to the
  true error term $\bm{\varepsilon}$, will be useful later on.

\item The OLS estimate of $\sigma^2$ (the variance of the error term), denoted
  $s^2$, is the sum of squared residuals divided by $n - K$:
  \begin{equation}
    s^2 \equiv \frac{SSR}{n - K} = \frac{\vec{e}'\vec{e}}{n - K}.
    \label{eq:1.2.13}
  \end{equation}
  (The definition presumes that $n > K$; otherwise $s^2$ is not well-defined.)
  As will be shown in Proposition~1.2 below, dividing the sum of squared residuals
  by $n - K$ (called the \textbf{degrees of freedom}) rather than by~$n$ (the
  sample size) makes this estimate unbiased for~$\sigma^2$.  The intuitive reason
  is that $K$ parameters ($\bm{\beta}$) have to be estimated before obtaining the
  residual vector~$\vec{e}$ used to calculate~$s^2$.  More specifically, $\vec{e}$
  has to satisfy the $K$ normal equations \eqref{eq:1.2.3prime}, which limits the
  variability of the residual.

\item The square root of $s^2$, $s$, is called the \textbf{standard error of the
  regression (SER)} or \textbf{standard error of the equation (SEE)}.  It is an
  estimate of the standard deviation of the error term.

\item The \textbf{sampling error} is defined as $\vec{b} - \bm{\beta}$.  It too
  can be related to $\bm{\varepsilon}$ as follows.
  \begin{align}
    \vec{b} - \bm{\beta}
      &= (\vec{X}'\vec{X})^{-1}\vec{X}'\vec{y} - \bm{\beta}
      &&\text{(by \eqref{eq:1.2.5})} \notag \\
      &= (\vec{X}'\vec{X})^{-1}\vec{X}'(\vec{X}\bm{\beta} + \bm{\varepsilon})
         - \bm{\beta}
      &&\text{(since $\vec{y} = \vec{X}\bm{\beta} + \bm{\varepsilon}$
         by Assumption~1.1)} \notag \\
      &= (\vec{X}'\vec{X})^{-1}(\vec{X}'\vec{X})\bm{\beta}
         + (\vec{X}'\vec{X})^{-1}\vec{X}'\bm{\varepsilon}
         - \bm{\beta} \notag \\
      &= \bm{\beta} + (\vec{X}'\vec{X})^{-1}\vec{X}'\bm{\varepsilon}
         - \bm{\beta}
       = (\vec{X}'\vec{X})^{-1}\vec{X}'\bm{\varepsilon}.
    \label{eq:1.2.14}
  \end{align}

\item \textbf{Uncentered $R^2$}.  One measure of the variability of the dependent
  variable is the sum of squares, $\sum y_i^2 = \vec{y}'\vec{y}$.  Because the OLS
  residual is chosen to satisfy the normal equations, we have the following
  decomposition of $\vec{y}'\vec{y}$:
  \begin{align}
    \vec{y}'\vec{y}
      &= (\hat{\vec{y}} + \vec{e})'(\hat{\vec{y}} + \vec{e})
      &&\text{(since $\vec{e} = \vec{y} - \hat{\vec{y}}$)} \notag \\
      &= \hat{\vec{y}}'\hat{\vec{y}} + 2\hat{\vec{y}}'\vec{e} + \vec{e}'\vec{e}
         \notag \\
      &= \hat{\vec{y}}'\hat{\vec{y}} + 2\vec{b}'\vec{X}'\vec{e} + \vec{e}'\vec{e}
      &&\text{(since $\hat{\vec{y}} \equiv \vec{X}\vec{b}$)} \notag \\
      &= \hat{\vec{y}}'\hat{\vec{y}} + \vec{e}'\vec{e}
      &&\text{(since $\vec{X}'\vec{e} = \vec{0}$ by the normal equations;
         see \eqref{eq:1.2.3prime})}.
    \label{eq:1.2.15}
  \end{align}

  The \textbf{uncentered $R^2$} is defined as
  \begin{equation}
    R_{uc}^2 \equiv 1 - \frac{\vec{e}'\vec{e}}{\vec{y}'\vec{y}}.
    \label{eq:1.2.16}
  \end{equation}
  Because of the decomposition \eqref{eq:1.2.15}, this equals
  \[
    \frac{\hat{\vec{y}}'\hat{\vec{y}}}{\vec{y}'\vec{y}}.
  \]
  Since both $\hat{\vec{y}}'\hat{\vec{y}}$ and $\vec{e}'\vec{e}$ are nonnegative,
  $0 \le R_{uc}^2 \le 1$.  Thus, the uncentered~$R^2$ has the interpretation of the
  fraction of the variation of the dependent variable that is attributable to the
  variation in the explanatory variables.  The closer the fitted value tracks the
  dependent variable, the closer the uncentered~$R^2$ to one.

\item \textbf{(Centered) $R^2$, the coefficient of determination}.  If the only
  regressor is a constant (so that $K = 1$ and $x_{i1} = 1$), then it is easy to
  see from \eqref{eq:1.2.5} that $\vec{b}$ equals $\bar{y}$, the sample mean of
  the dependent variable, which means that $\hat{y}_i = \bar{y}$ for all~$i$,
  $\hat{\vec{y}}'\hat{\vec{y}}$ in \eqref{eq:1.2.15} equals $n\bar{y}^2$, and
  $\vec{e}'\vec{e}$ equals $\sum_i (y_i - \bar{y})^2$.  If the regressors also
  include nonconstant variables, then it can be shown (the proof is left as an
  analytical exercise) that $\sum_i (y_i - \bar{y})^2$ is decomposed as
  \begin{equation}
    \sum_{i=1}^{n}(y_i - \bar{y})^2
      = \sum_{i=1}^{n}(\hat{y}_i - \bar{y})^2
        + \sum_{i=1}^{n} e_i^2
    \quad\text{with}\quad
    \bar{y} \equiv \frac{1}{n}\sum_{i=1}^{n} y_i.
    \label{eq:1.2.17}
  \end{equation}

  The \textbf{coefficient of determination}, $R^2$, is defined as
  \begin{equation}
    R^2 \equiv 1 - \frac{\sum_{i=1}^{n} e_i^2}
                        {\sum_{i=1}^{n}(y_i - \bar{y})^2}.
    \label{eq:1.2.18}
  \end{equation}
  Because of the decomposition \eqref{eq:1.2.17}, this $R^2$ equals
  \[
    \frac{\sum_{i=1}^{n}(\hat{y}_i - \bar{y})^2}
         {\sum_{i=1}^{n}(y_i - \bar{y})^2}.
  \]
  Therefore, provided that the regressors include a constant so that the
  decomposition \eqref{eq:1.2.17} is valid, $0 \le R^2 \le 1$.  Thus, this $R^2$
  as defined in \eqref{eq:1.2.18} is a measure of the explanatory power of the
  nonconstant regressors.

  If the regressors do not include a constant but (as some regression software
  packages do) you nevertheless calculate $R^2$ by the formula \eqref{eq:1.2.18},
  then the $R^2$ can be negative.  This is because, without the benefit of an
  intercept, the regression could do worse than the sample mean in terms of
  tracking the dependent variable.  On the other hand, some other regression
  packages (notably STATA) switch to the formula \eqref{eq:1.2.16} for the $R^2$
  when a constant is not included, in order to avoid negative values for the
  $R^2$.  This is a mixed blessing.  Suppose that the regressors do not include a
  constant but that a linear combination of the regressors equals a constant.
  This occurs if, for example, the intercept is replaced by seasonal
  dummies.\footnote{Dummy variables will be introduced in the empirical exercise
  for this chapter.}  The regression is essentially the same when one of the
  regressors in the linear combination is replaced by a constant.  Indeed, one
  should obtain the same vector of fitted values.  But if the formula for the
  $R^2$ is \eqref{eq:1.2.16} for regressions without a constant and
  \eqref{eq:1.2.18} for those with a constant, the calculated $R^2$ declines
  (see Review Question~7 below) after the replacement by a constant.

\end{itemize}

%% ─── Influential Analysis (optional) ─────────────────────────────────────────
\subsection*{Influential Analysis (optional)}

Since the method of least squares seeks to prevent a few large residuals at the
expense of incurring many relatively small residuals, only a few observations can
be extremely influential in the sense that dropping them from the sample changes
some elements of $\vec{b}$ substantially.  There is a systematic way to find those
\textbf{influential observations}.\footnote{See Krasker, Kuh, and Welsch (1983)
for more details.}  Let $\vec{b}^{(i)}$ be the OLS estimate of $\bm{\beta}$ that
would be obtained if OLS were used on a sample from which the $i$-th observation
was omitted.  The key equation is
\begin{equation}
  \vec{b}^{(i)} - \vec{b}
    = -\left(\frac{1}{1 - p_i}\right)
      (\vec{X}'\vec{X})^{-1}\vec{x}_i \cdot e_i,
  \label{eq:1.2.19}
\end{equation}
where $\vec{x}_i$ as before is the $i$-th row of $\vec{X}$, $e_i$ is the OLS
residual for observation~$i$, and $p_i$ is defined as
\begin{equation}
  p_i \equiv \vec{x}_i'(\vec{X}'\vec{X})^{-1}\vec{x}_i,
  \label{eq:1.2.20}
\end{equation}
which is the $i$-th diagonal element of the projection matrix~$\vec{P}$.
(Proving \eqref{eq:1.2.19} would be a good exercise in matrix algebra, but we
will not do it here.)  It is easy to show (see Review Question~7 of Section~1.3)
that
\begin{equation}
  0 \le p_i \le 1
  \quad\text{and}\quad
  \sum_{i=1}^{n} p_i = K.
  \label{eq:1.2.21}
\end{equation}
So $p_i$ equals $K/n$ on average.

To illustrate the use of \eqref{eq:1.2.19} in a specific example, consider the
relationship between equipment investment and economic growth for the world's
poorest countries between 1960 and 1985.  Figure~\ref{fig:1.2} plots the average
annual GDP-per-worker growth between 1960 and 1985 against the ratio of equipment
investment to GDP over the same period for thirteen countries whose GDP per worker
in 1965 was less than 10 percent of that of the United
States.\footnote{The data are from the Penn World Table, reprinted in DeLong and
Summers (1991).  To their credit, their analysis is based on the whole sample of
sixty-one countries.}  It is clear visually from the plot that the position of the
estimated regression line would depend very much on the single outlier (Botswana).
Indeed, if Botswana is dropped from the sample, the estimated slope coefficient
drops from 0.37 to 0.058.  In the present case of simple regression, it is easy to
spot outliers by visually inspecting the plot such as Figure~\ref{fig:1.2}.  This
strategy would not work if there were more than one nonconstant regressor.
Analysis based on formula \eqref{eq:1.2.19} is not restricted to simple
regressions.  Table~\ref{tab:1.1} displays the data along with the OLS residuals,
the values of~$p_i$, and \eqref{eq:1.2.19} for each observation.  Botswana's
$p_i$ of 0.7196 is well above the average of $0.154$ ($= K/n = 2/13$) and is
highly \textbf{influential}, as the last two columns of the table indicate.  Note
that we could not have detected the influential observation by looking at the
residuals, which is not surprising because the algebra of least squares is designed
to avoid large residuals at the expense of many small residuals for other
observations.

What should be done with influential observations?  It depends.  If the
influential observations satisfy the regression model, they provide valuable
information about the regression function unavailable from the rest of the sample
and should definitely be kept in the sample.  But more probable is that the
influential observations are atypical of the rest of the sample because they do
not satisfy the model.  In this case they should definitely be dropped from the
sample.  For the example just examined, there was a worldwide growth in the demand
for diamonds, Botswana's main export, and production of diamonds requires heavy
investment in drilling equipment.  If the reason to expect an association between
growth and equipment investment is the beneficial effect on productivity of the
introduction of new technologies through equipment, then Botswana, whose high GDP
growth is demand-driven, should be dropped from the sample.

\begin{figure}[htbp]
\centering
\includegraphics[width=0.9\textwidth,trim=50 400 50 400,clip]{figures/fig_1_2.png}
\caption{Equipment Investment and Growth}
\label{fig:1.2}
\end{figure}

\begin{table}[htbp]
\centering
\caption{Influential Analysis}
\label{tab:1.1}
\small
\begin{tabular}{l r r r r r r}
\toprule
  & GDP/worker & Equipment/ & & & (1.2.19) & (1.2.19) \\
Country & growth & GDP & Residual & $p_i$ & for $\beta_1$ & for $\beta_2$ \\
\midrule
Botswana    &  0.0676 & 0.1310 &  0.0119 & 0.7196 &  0.0104 & $-0.3124$ \\
Cameroon    &  0.0458 & 0.0415 &  0.0233 & 0.0773 & $-0.0021$ &  0.0045 \\
Ethiopia    &  0.0094 & 0.0212 & $-0.0056$ & 0.1193 &  0.0010 & $-0.0119$ \\
India       &  0.0115 & 0.0278 & $-0.0059$ & 0.0980 &  0.0009 & $-0.0087$ \\
Indonesia   &  0.0345 & 0.0221 &  0.0192 & 0.1160 & $-0.0034$ &  0.0394 \\
Ivory Coast &  0.0278 & 0.0243 &  0.0117 & 0.1084 & $-0.0019$ &  0.0213 \\
Kenya       &  0.0146 & 0.0462 & $-0.0096$ & 0.0775 &  0.0007 &  0.0023 \\
Madagascar  & $-0.0102$ & 0.0219 & $-0.0254$ & 0.1167 &  0.0045 & $-0.0527$ \\
Malawi      &  0.0153 & 0.0361 & $-0.0052$ & 0.0817 &  0.0006 & $-0.0036$ \\
Mali        &  0.0044 & 0.0433 & $-0.0188$ & 0.0769 &  0.0016 & $-0.0006$ \\
Pakistan    &  0.0295 & 0.0263 &  0.0126 & 0.1022 & $-0.0020$ &  0.0205 \\
Tanzania    &  0.0184 & 0.0860 & $-0.0206$ & 0.2281 & $-0.0021$ &  0.0952 \\
Thailand    &  0.0341 & 0.0395 &  0.0123 & 0.0784 & $-0.0012$ &  0.0047 \\
\bottomrule
\end{tabular}
\end{table}

%% ─── A Note on the Computation of OLS Estimates ──────────────────────────────
\subsection*{A Note on the Computation of OLS Estimates}

So far, we have focused on the conceptual aspects of the algebra of least squares.
But for applied researchers who actually calculate OLS estimates using digital
computers, it is important to be aware of a certain aspect of digital computing in
order to avoid the risk of obtaining unreliable estimates without knowing it.  The
source of a potential problem is that the computer approximates real numbers by
so-called \textbf{floating-point numbers}.  When an arithmetic operation involves
both very large numbers and very small numbers, floating-point calculation can
produce inaccurate results.  This is relevant in the computation of OLS estimates
when the regressors greatly differ in magnitude.  For example, one of the
regressors may be the interest rate stated as a fraction, and another may be U.S.\
GDP in dollars.  The matrix $\vec{X}'\vec{X}$ will then contain both very small
and very large numbers, and the arithmetic operation of inverting this matrix by
the digital computer will produce unreliable results.

A simple solution to this problem is to choose the units of measurement so that
the regressors are similar in magnitude.  For example, state the interest rate in
percents and U.S.\ GDP in trillion dollars.  This sort of care would prevent the
problem most of the time.  A more systematic transformation of the $\vec{X}$
matrix is to subtract the sample means of all regressors and divide by the sample
standard deviations before forming $\vec{X}'\vec{X}$ (and adjust the OLS estimates
to undo the transformation).  Most OLS programs (such as TSP) take a more
sophisticated transformation of the $\vec{X}$ matrix (called the
\textbf{QR decomposition}) to produce accurate results.
